\documentclass[a4paper,11pt]{article}
\usepackage[utf8]{inputenc}
\usepackage[T1]{fontenc}
\usepackage[french]{babel}
\usepackage[left=2cm,right=2cm,top=2cm,bottom=2cm]{geometry}
\usepackage{xcolor}
\usepackage{hyperref}
\usepackage{fontawesome5}
\usepackage{parskip}

% --- Couleurs Naval Group ---
\definecolor{primary}{RGB}{0, 51, 102}
\hypersetup{colorlinks=true, urlcolor=primary}

\begin{document}

% --- EXPÉDITEUR ---
\noindent
\begin{minipage}[t]{0.5\textwidth}
    \textbf{Loïc Aron MBASSI EWOLO}\\
    Élève-Ingénieur Bac+5 -- Double diplôme ENSEA\\[3pt]
    \faMapMarker\ Cergy, France\\
    \faPhone\ (+33) 7 59 52 33 98\\
    \faEnvelope\ \href{mailto:loicaron.mbassiewolo@ensea.fr}{loicaron.mbassiewolo@ensea.fr}
\end{minipage}
\hfill
% --- DATE ---
\begin{minipage}[t]{0.4\textwidth}
    \raggedleft
    Cergy, le 5 février 2026
\end{minipage}

\vspace{1.2cm}

% --- DESTINATAIRE ---
\hfill
\begin{minipage}[t]{0.5\textwidth}
    \raggedleft
    \textbf{Naval Group}\\
    Service Recrutement\\
    Ollioules, France
\end{minipage}

\vspace{1cm}

\noindent\textbf{Objet :} Candidature au stage Développeur Logiciel Embarqué -- Guerre des Mines

\vspace{0.8cm}

Madame, Monsieur,

Élève-ingénieur en double diplôme à l'ENSEA et à l'École Polytechnique de Yaoundé, spécialisé en systèmes embarqués et développement logiciel, je souhaite rejoindre le service DILM de Naval Group à Ollioules pour contribuer au développement des systèmes de mission dans le domaine de la guerre des mines. La perspective de travailler sur des systèmes interagissant avec des drones multi-milieux et des systèmes de combat représente un défi technique passionnant.

Mon expérience actuelle sur le challenge CoHoMa, en collaboration avec l'Armée Française et l'ENSEA, m'a immergé dans la robotique autonome de défense. J'y développe des solutions de cartographie, navigation autonome et fusion de capteurs (Lidar 3D, caméras de profondeur) en C/C++, avec optimisation IA embarquée sur GPU Nvidia Jetson. Cette expérience m'a confronté aux exigences des systèmes critiques militaires et au travail en équipe pluridisciplinaire.

Côté développement logiciel, je maîtrise Java, TypeScript et Angular, technologies centrales pour votre projet. J'ai conçu des architectures backend robustes avec Spring Boot et NestJS, en appliquant les design patterns (Singleton, Factory, Observer, MVC) dans une approche orientée objet rigoureuse. Mon projet d'optimisation de routes de collecte, récompensé du 1\textsuperscript{er} prix national CITS, démontre ma capacité à concevoir des algorithmes complexes intégrables dans des systèmes opérationnels, compétence directement applicable à l'optimisation de la chasse aux mines.

Je suis particulièrement motivé par les deux sujets proposés : l'optimisation du découpage de zones via cartographie ou langage naturel, et l'exploitation de la bathymétrie et des courants marins pour le minage. Mon expérience en IA (PyTorch, NLP) et mon système multi-agents orchestré avec RAG me permettraient de contribuer efficacement à ces fonctionnalités innovantes. La méthodologie Agile avec sprints de 3 semaines correspond parfaitement à mon expérience de travail en équipe projet.

Disponible dès avril 2026 pour 4 à 6 mois, je serais honoré d'intégrer l'équipe de Michael et de contribuer aux systèmes de mission de Naval Group.

Je vous prie d'agréer, Madame, Monsieur, l'expression de mes salutations distinguées.

\vspace{0.8cm}

\textbf{Loïc Aron MBASSI EWOLO}\\
\textit{Élève-Ingénieur ENSEA / Polytechnique Yaoundé}

\end{document}
